\chapter{进程同步}

\section{同步与互斥的概念}
\concept{并发进程的制约关系}
\begin{points}
  \item 多道程序系统中，由于资源共享或进程合作，并发进程之间形成相互制约关系。
  \item 间接相互制约：多个进程竞争同一资源，例如多个进程争用打印机。
  \item 直接相互制约：进程之间存在合作顺序，例如生产者生产后消费者才能消费。
  \item 进程同步的任务是协调这些关系，使并发进程能正确共享资源和相互合作，保证结果可再现。
\end{points}
\plain{核心是“看起来同时推进”。单 CPU 上靠时间片交替，多 CPU 上才可能真的同一时刻运行。}
\exam{常考并发与并行区别、单处理机下最多几个进程运行、并发导致为什么需要同步互斥。}


\concept{与时间有关的错误/race condition}
\begin{points}
  \item 多个进程并发读写共享变量，执行次序不同可能导致不同结果，这种错误称为与时间有关的错误或竞争条件。
  \item 例子：订票系统中 A、B 同时读取剩余票数 $x=1$，都判断有票并出票，最终可能超卖。
  \item 根本原因：对共享变量的复合操作没有互斥保护。
  \item 解决思路：把读-改-写共享变量的代码段放入临界区，一次只允许一个进程进入。
\end{points}
\plain{并发程序共享数据时，执行顺序一变结果就变，所以要把关键代码段保护起来。}
\exam{常考直接/间接制约、竞争条件、临界资源/临界区、四条临界区准则。}


\concept{临界资源与临界区}
\begin{points}
  \item 临界资源是一段时间内只允许一个进程使用的资源。
  \item 临界区是进程中访问临界资源的代码段。
  \item 进入区负责申请进入临界区；退出区负责释放临界区；剩余区是其他代码。
  \item 互斥要求：任意时刻最多一个进程在相关临界区内执行。
\end{points}
\plain{并发程序共享数据时，执行顺序一变结果就变，所以要把关键代码段保护起来。}
\exam{常考直接/间接制约、竞争条件、临界资源/临界区、四条临界区准则。}


\concept{临界区访问原则}
\begin{points}
  \item \key{空闲让进}：临界区空闲时，应允许一个请求进入的进程立即进入。
  \item \key{忙则等待}：临界区已有进程时，其他请求进程必须等待。
  \item \key{有限等待}：等待进入临界区的进程不能无限期等待，避免饥饿。
  \item \key{让权等待}：不能进入临界区的进程应释放处理机，不应长时间忙等浪费 CPU。
\end{points}
\plain{并发程序共享数据时，执行顺序一变结果就变，所以要把关键代码段保护起来。}
\exam{常考直接/间接制约、竞争条件、临界资源/临界区、四条临界区准则。}

\section{互斥的实现方法}
\concept{软件方法}
\begin{points}
  \item 单标志法、双标志先检查法、双标志后检查法、Peterson 算法都是早期软件互斥方法。
  \item 软件方法的关键是用共享变量表达“谁想进”和“轮到谁进”。
  \item Peterson 算法能满足互斥、空闲让进和有限等待，但主要适合两个进程的理论模型。
  \item 软件方法缺点是复杂、容易出错、依赖忙等，难以推广。
\end{points}
\plain{互斥可以靠软件协议，也可以靠关中断/测试并设置等硬件原子指令；现代更常封装成锁或信号量。}
\exam{常考忙等缺点、关中断是否适合多处理器、Test-and-Set 的作用。}


\concept{硬件方法}
\begin{points}
  \item 关中断：单处理器中进入临界区前关中断，退出后开中断，可防止进程被中断切换。
  \item Test-and-Set 指令：原子地测试并设置锁变量。
  \item Swap 指令：原子地交换两个变量的值。
  \item 硬件方法可保证互斥，但若采用自旋锁，会造成忙等，不满足让权等待。
\end{points}
\plain{互斥可以靠软件协议，也可以靠关中断/测试并设置等硬件原子指令；现代更常封装成锁或信号量。}
\exam{常考忙等缺点、关中断是否适合多处理器、Test-and-Set 的作用。}


\section{信号量与 P/V 操作}
\concept{信号量定义}
\begin{points}
  \item 信号量 semaphore 由整数值 $s$ 和等待队列 $q$ 组成。
  \item 除初始化外，信号量的值只能通过 P 操作和 V 操作改变。
  \item 信号量 $s>0$ 时，表示可用资源数；$s=0$ 表示无可用资源但无人等待；$s<0$ 时，$|s|$ 表示等待队列中的进程数。
  \item P/V 操作必须是原子操作，执行过程中不能被其他进程打断。
\end{points}
\plain{P 是先占/先减，不够就睡；V 是释放/先加，如果有人等就叫醒一个。}
\exam{常考 P/V 原语含义、s>0/s<0 的解释、互斥信号量和资源信号量的初值。}


\concept{P 操作}
\begin{points}
  \item P(S)：先执行 $s=s-1$。
  \item 若减完后 $s<0$，说明资源不够，当前进程阻塞，进入该信号量等待队列，并转调度程序。
  \item 若减完后 $s\ge 0$，说明成功获得资源，进程继续执行。
  \item 直觉：P 是“申请/占用一个资源”，先登记需求，再判断是否需要等待。
\end{points}
\plain{P 是先占/先减，不够就睡；V 是释放/先加，如果有人等就叫醒一个。}
\exam{常考 P/V 原语含义、s>0/s<0 的解释、互斥信号量和资源信号量的初值。}


\concept{V 操作}
\begin{points}
  \item V(S)：先执行 $s=s+1$。
  \item 若加完后 $s\le 0$，说明仍有进程在等待队列中，唤醒一个等待进程并插入就绪队列。
  \item 若加完后 $s>0$，说明没有进程等待，只是增加了一个可用资源。
  \item 直觉：V 是“释放一个资源”。即使加完后 $s\le 0$，也要唤醒，因为负数表示有等待者；释放的资源应交给其中一个等待进程。
\end{points}
\plain{P 是先占/先减，不够就睡；V 是释放/先加，如果有人等就叫醒一个。}
\exam{常考 P/V 原语含义、s>0/s<0 的解释、互斥信号量和资源信号量的初值。}

\concept{整型信号量与记录型信号量}
\begin{points}
  \item 整型信号量只用整数表示资源数，P 操作中可能忙等。
  \item 记录型信号量增加等待队列，资源不足时让进程阻塞，满足让权等待。
  \item 考试中常默认 P/V 是记录型信号量，阻塞和唤醒由 OS 完成。
\end{points}
\plain{P 是先占/先减，不够就睡；V 是释放/先加，如果有人等就叫醒一个。}
\exam{常考 P/V 原语含义、s>0/s<0 的解释、互斥信号量和资源信号量的初值。}


\section{经典同步问题}
\concept{生产者-消费者问题}
\begin{points}
  \item 问题：多个生产者向有界缓冲区放产品，多个消费者从缓冲区取产品。
  \item 需要三个信号量：\code{mutex=1} 保护缓冲区互斥访问，\code{empty=n} 表示空缓冲区数，\code{full=0} 表示满缓冲区数。
  \item 生产者顺序：生产产品 $\rightarrow$ P(empty) $\rightarrow$ P(mutex) $\rightarrow$ 放入缓冲区 $\rightarrow$ V(mutex) $\rightarrow$ V(full)。
  \item 消费者顺序：P(full) $\rightarrow$ P(mutex) $\rightarrow$ 取出产品 $\rightarrow$ V(mutex) $\rightarrow$ V(empty) $\rightarrow$ 消费产品。
  \item 注意：先申请同步资源 empty/full，再申请互斥 mutex，避免拿着锁等待条件导致其他进程无法改变条件。
\end{points}
\plain{经典同步题先找资源数量，再找互斥，再找先后顺序。}
\exam{常考补全 P/V 代码、判断会不会死锁、修改信号量初值。}


\concept{读者-写者问题}
\begin{points}
  \item 多个读者可同时读共享文件，但写者必须互斥写；写者与读者也互斥。
  \item 读者优先方案中，第一个读者负责阻止写者，最后一个读者负责释放写者。
  \item 需要计数变量 \code{readcount} 记录当前读者数，并用互斥信号量保护它。
  \item 易错点：读者之间不是互斥读文件，而是互斥修改 \code{readcount}。
\end{points}
\plain{经典同步题先找资源数量，再找互斥，再找先后顺序。}
\exam{常考补全 P/V 代码、判断会不会死锁、修改信号量初值。}


\concept{哲学家进餐问题}
\begin{points}
  \item 五个哲学家围桌，每人左右各一只筷子，必须同时拿到两只筷子才能进餐。
  \item 若所有哲学家都先拿左筷子再等右筷子，可能形成死锁。
  \item 解决方法包括：最多允许四人同时取筷子；奇偶哲学家取筷子顺序相反；一次性申请两只筷子；使用管程。
  \item 核心直觉：避免每个进程“占有一个资源，同时等待另一个资源”。
\end{points}
\plain{经典同步题先找资源数量，再找互斥，再找先后顺序。}
\exam{常考补全 P/V 代码、判断会不会死锁、修改信号量初值。}


\section{信号量集机制}
\concept{AND 型信号量}
\begin{points}
  \item AND 型信号量适用于进程同时需要多类资源、每类一个的情况。
  \item 基本思想：进程需要的多类资源一次性全部分配；只要有一个资源不满足，其他资源也不分配。
  \item 这样可以避免进程占有部分资源再等待其他资源，从而降低死锁风险。
  \item P 原语称为 SP/Swait，V 原语称为 SV/Ssignal。
\end{points}
\plain{P 是先占/先减，不够就睡；V 是释放/先加，如果有人等就叫醒一个。}
\exam{常考 P/V 原语含义、s>0/s<0 的解释、互斥信号量和资源信号量的初值。}


\concept{一般信号量集}
\begin{points}
  \item 一般信号量集是 AND 型的扩展：一次可申请多类资源，每类资源可申请多个。
  \item 当某类资源低于下限或不能满足申请量时，本次原语不分配任何资源。
  \item 它把“多个 P 操作”合并成一个原子判断，避免中间状态。
\end{points}
\plain{P 是先占/先减，不够就睡；V 是释放/先加，如果有人等就叫醒一个。}
\exam{常考 P/V 原语含义、s>0/s<0 的解释、互斥信号量和资源信号量的初值。}


\section{管程}
\concept{管程是什么}
\begin{points}
  \item 管程是一种高级同步机制，把共享数据、对共享数据的操作、同步条件封装在一个模块中。
  \item 管程保证任意时刻只有一个进程/线程在管程内部执行，因此天然提供互斥。
  \item 条件变量用于表达“某个条件不满足时等待，条件满足时被唤醒”。
  \item 典型操作包括 \code{wait} 和 \code{signal}。
\end{points}
\plain{管程把共享数据和操作封装起来，自动保证一次只有一个进程进管程。}
\exam{常考管程和 P/V 区别：管程是高级语言级封装，P/V 更底层更灵活也更易错。}


\concept{管程与 P/V 的区别}
\begin{points}
  \item P/V 是低级同步原语，程序员需要自己安排每个 P、V 的位置，容易写错。
  \item 管程把共享变量和互斥规则封装起来，进入管程自动互斥，程序员主要表达条件等待。
  \item P/V 更像“手动锁门开门”；管程更像“把房间设计成一次只能进一个人，并提供排队等待条件”。
  \item 管程更结构化，适合高级语言；信号量更基础，适合描述底层同步。
\end{points}
\plain{管程把共享数据和操作封装起来，自动保证一次只有一个进程进管程。}
\exam{常考管程和 P/V 区别：管程是高级语言级封装，P/V 更底层更灵活也更易错。}


\chapter{死锁}
